\unitchapter{2}{7}{Data Structures \& Algorithms}{p2-07}

\unitmeta{Expected questions: 10--12 · Target: 9+ · Type: complexity, recurrences,
tree/heap/hash numericals, algorithm traces}

\begin{sylbox}

\begin{itemize}
\tightlist
\item[$\square$]
  Data structures: arrays, sparse matrices, stacks, queues, priority queues, linked lists, trees, binary trees, threaded trees, BST, AVL, B/B+/B* trees, sets (union-find), graphs, sorting, searching, hashing
\item[$\square$]
  Performance analysis: time \& space complexity, asymptotic notation, recurrence relations
\item[$\square$]
  Design techniques: divide \& conquer, dynamic programming, greedy, backtracking, branch \& bound
\item[$\square$]
  Lower bound theory: comparison trees, reductions
\item[$\square$]
  Graph algorithms: BFS, DFS, shortest paths, maximum flow, MST
\item[$\square$]
  Complexity theory: P, NP, NP-complete, NP-hard, reducibility
\item[$\square$]
  Selected topics: number-theoretic algorithms, polynomial arithmetic, FFT, string matching
\item[$\square$]
  Advanced: parallel algorithms (sorting, searching, merging), approximation \& randomised algorithms
\end{itemize}

\end{sylbox}

\needspace{125pt}

\hypertarget{p2-07--1-asymptotic-analysis}{%
\section{Asymptotic Analysis}\label{p2-07--1-asymptotic-analysis}}

Asymptotic analysis describes how an algorithm\textquotesingle s running time grows with the input size \(n\), ignoring constant
factors and small inputs. It lets us compare algorithms independently of the machine they run on.

\begin{formulatable}
$O$ --- upper bound & f(n) \le c\,g(n) \quad\text{for all } n \ge n_0 \\
$\Omega$ --- lower bound & f(n) \ge c\,g(n) \quad\text{for all } n \ge n_0 \\
$\Theta$ --- tight bound & c_1 g(n) \le f(n) \le c_2 g(n) \\
$o$ --- strict upper bound & \lim_{n\to\infty} f(n)/g(n) = 0 \\
$\omega$ --- strict lower bound & \lim_{n\to\infty} f(n)/g(n) = \infty \\
\end{formulatable}


Growth order (memorise):

\[
1 < \log\log n < \log n < (\log n)^k < \sqrt n < n < n\log n < n^2 < n^3 < 2^n < 3^n < n! < n^n
\]
\begin{figure}[!htbp]
\centering
\begin{adjustbox}{max width=\linewidth}
\begin{tikzpicture}
  \begin{axis}[width=10.5cm, height=6.4cm, xmin=1, xmax=16, ymin=0, ymax=70, axis lines=left,
      xlabel={$n$}, ylabel={operations}, tick label style={font=\footnotesize}, label style={font=\small},
      legend style={font=\footnotesize, draw=none, fill=none, at={(1.02,1)}, anchor=north west}, legend cell align=left,
      every axis plot/.append style={line width=0.9pt, no markers, samples=80, domain=1:16}]
    \addplot[fgink!55] {log2(x)}; \addlegendentry{$\log n$}
    \addplot[fsteel!80!black] {x}; \addlegendentry{$n$}
    \addplot[fmoss!80!black] {x*log2(x)}; \addlegendentry{$n\log n$}
    \addplot[fochre!90!black, domain=1:8.4] {x^2}; \addlegendentry{$n^2$}
    \addplot[fwine, domain=1:6.13] {2^x}; \addlegendentry{$2^n$}
  \end{axis}
\end{tikzpicture}
\end{adjustbox}
\caption{Growth of common running times. For small $n$ the curves cross; asymptotic order describes only
what happens as $n$ grows large.}
\end{figure}
\noindent Useful facts: $\log(n!) = \Theta(n\log n)$; \ $n^{1/\log n} = \Theta(1)$; \ $2^{\log_2 n} = n$; \
$n^{\log n}$ is super-polynomial.


\needspace{160pt}

\hypertarget{p2-07--recurrences--master-theorem}{%
\subsection{Recurrences --- Master Theorem}\label{p2-07--recurrences--master-theorem}}

\(T(n) = a\,T(n/b) + f(n)\), \(a \ge 1\), \(b > 1\). Compare \(f(n)\) with \textbf{\(n^{\log_b a}\)}:

\begingroup\small

\par\medskip\noindent\hfil\begin{tabular}{@{}
  >{\raggedright\arraybackslash}p{(\columnwidth - 4\tabcolsep) * \real{0.0655}}
  >{\raggedright\arraybackslash}p{(\columnwidth - 4\tabcolsep) * \real{0.5080}}
  >{\raggedright\arraybackslash}p{(\columnwidth - 4\tabcolsep) * \real{0.3155}}@{}}
\toprule\noalign{}
\begin{minipage}[b]{\linewidth}\raggedright
\textbf{Case}
\end{minipage} & \begin{minipage}[b]{\linewidth}\raggedright
\textbf{Condition}
\end{minipage} & \begin{minipage}[b]{\linewidth}\raggedright
\textbf{Result}
\end{minipage} \\
\midrule\noalign{}
1 & \(f(n) = O(n^{\log_b a - \varepsilon})\) & \(\Theta(n^{\log_b a})\) \\
2 & \(f(n) = \Theta(n^{\log_b a} \log^k n)\) & \(\Theta(n^{\log_b a} \log^{k+1} n)\) \\
3 & \(f(n) = \Omega(n^{\log_b a + \varepsilon})\) and regularity & \(\Theta(f(n))\) \\
\bottomrule\noalign{}
\end{tabular}\hfil\null\par\medskip


\endgroup

\begingroup\small

\par\medskip\noindent\hfil\begin{tabular}{@{}
  >{\raggedright\arraybackslash}p{(\columnwidth - 4\tabcolsep) * \real{0.2741}}
  >{\raggedright\arraybackslash}p{(\columnwidth - 4\tabcolsep) * \real{0.2948}}
  >{\raggedright\arraybackslash}p{(\columnwidth - 4\tabcolsep) * \real{0.2575}}@{}}
\toprule\noalign{}
\begin{minipage}[b]{\linewidth}\raggedright
\textbf{Recurrence}
\end{minipage} & \begin{minipage}[b]{\linewidth}\raggedright
\textbf{Solution}
\end{minipage} & \begin{minipage}[b]{\linewidth}\raggedright
\textbf{Algorithm}
\end{minipage} \\
\midrule\noalign{}
\(T(n) = T(n/2) + 1\) & \(\Theta(\log n)\) & Binary search \\
\(T(n) = 2T(n/2) + n\) & \(\Theta(n \log n)\) & Merge sort \\
\(T(n) = 2T(n/2) + 1\) & \(\Theta(n)\) & Tree traversal, max-min \\
\(T(n) = T(n-1) + n\) & \(\Theta(n^2)\) & Quick sort worst, selection \\
\(T(n) = T(n-1) + 1\) & \(\Theta(n)\) & Linear recursion \\
\(T(n) = 2T(n-1) + 1\) & \(\Theta(2^n)\) & Tower of Hanoi \\
\(T(n) = 7T(n/2) + n^2\) & \(\Theta(n^{2.81})\) & Strassen \\
\(T(n) = 8T(n/2) + n^2\) & \(\Theta(n^3)\) & Naive D\&C matrix multiply \\
\(T(n) = 3T(n/2) + n\) & \(\Theta(n^{1.585})\) & Karatsuba \\
\(T(n) = T(\sqrt{n}) + 1\) & \(\Theta(\log \log n)\) & --- \\
\(T(n) = 2T(\sqrt{n}) + \log n\) & \(\Theta(\log n \cdot \log \log n)\) & --- \\
\(T(n) = T(n/3) + T(2n/3) + n\) & \(\Theta(n \log n)\) & Recursion tree \\
\bottomrule\noalign{}
\end{tabular}\hfil\null\par\medskip


\endgroup

\needspace{161pt}

\hypertarget{p2-07--2-linear-data-structures}{%
\section{Linear Data Structures}\label{p2-07--2-linear-data-structures}}

Linear structures arrange items in a sequence. They differ in where items may be inserted and removed: anywhere
(array, list), at one end only (stack) or at opposite ends (queue).

\needspace{95pt}

\hypertarget{p2-07--21-arrays}{%
\subsection{Arrays}\label{p2-07--21-arrays}}

\begin{itemize}
\tightlist
\item
  Row-major address: \textbf{A{[}i{]}{[}j{]} = B + w·{[}(i − L₁)·N + (j − L₂){]}} where N = number of columns.
\item
  Column-major: \textbf{B + w·{[}(j − L₂)·M + (i − L₁){]}}, M = rows.
\item
  Lower triangular n×n stored row-wise: index of (i, j) = i(i−1)/2 + j (1-based). Total elements = n(n+1)/2.
\item
  \textbf{Sparse matrix}: triplet (row, col, value) representation; or linked lists.
\end{itemize}

\needspace{95pt}

\hypertarget{p2-07--22-stack-lifo}{%
\subsection{Stack (LIFO)}\label{p2-07--22-stack-lifo}}

\begin{itemize}
\tightlist
\item
  Applications: expression conversion/evaluation, recursion, backtracking, parenthesis matching, undo.
\item
  \textbf{Infix → Postfix} example: \texttt{A\ +\ B\ *\ C\ −\ D\ /\ E} → \textbf{\texttt{A\ B\ C\ *\ +\ D\ E\ /\ −}}.
\item
  Postfix evaluation \texttt{6\ 2\ 3\ +\ −\ 3\ 8\ 2\ /\ +\ *} → 6 − 5 = 1; 8/2 = 4; 3 + 4 = 7; 1 × 7 = \textbf{7}.
\item
  Number of stack permutations of n elements = Catalan number Cₙ.
\end{itemize}

\needspace{95pt}

\hypertarget{p2-07--23-queue-fifo}{%
\subsection{Queue (FIFO)}\label{p2-07--23-queue-fifo}}

\begin{itemize}
\tightlist
\item
  Circular queue: full when (rear + 1) \% n == front (one slot wasted). Deque, priority queue.
\item
  Queue using two stacks: amortised O(1).
\end{itemize}

\needspace{227pt}

\hypertarget{p2-07--24-linked-lists}{%
\subsection{Linked Lists}\label{p2-07--24-linked-lists}}

Singly, doubly, circular. Insert at head O(1); search O(n). Reverse a list: three pointers. Detect cycle: Floyd\textquotesingle s tortoise-hare. Middle: slow/fast pointers.

\needspace{161pt}

\hypertarget{p2-07--3-trees}{%
\section{Trees}\label{p2-07--3-trees}}

A tree stores items hierarchically: one root, and every other node has exactly one parent. Trees give
logarithmic search when balanced, and they underlie heaps, indexes and syntax trees.

\needspace{95pt}

\hypertarget{p2-07--31-binary-tree-facts}{%
\subsection{Binary Tree Facts}\label{p2-07--31-binary-tree-facts}}

\begin{formulatable}
Maximum nodes at level $i$ (root at level 0) & 2^{i} \\
Maximum nodes in a tree of height $h$ (root height 0) & 2^{h+1} - 1 \\
Minimum height with $n$ nodes & \lceil \log_2(n+1) \rceil - 1 \\
Leaves versus nodes with two children & n_0 = n_2 + 1 \\
Null pointers in an $n$-node binary tree & n + 1 \\
Unlabelled binary trees with $n$ nodes (Catalan number) & C_n = \frac{1}{n+1}\binom{2n}{n} \\
Labelled binary trees with $n$ nodes & n!\cdot C_n \\
Full $k$-ary tree with $i$ internal nodes & n = k\,i + 1,\quad L = (k-1)\,i + 1 \\
\end{formulatable}


\needspace{95pt}

\hypertarget{p2-07--32-traversals}{%
\subsection{Traversals}\label{p2-07--32-traversals}}

\begin{figure}[!htbp]
\centering
\begin{adjustbox}{max width=\linewidth}
\begin{tikzpicture}[level distance=13mm, tn/.style={tnode, fill=fslate},
    level 1/.style={sibling distance=30mm}, level 2/.style={sibling distance=16mm}, edge from parent/.style={draw=fgink, line width=0.55pt}]
  \node[tn] (A) {A}
    child { node[tn] (B) {B} child { node[tn] (D) {D} } child { node[tn] (E) {E} } }
    child { node[tn] (C) {C} child[missing] child { node[tn] (F) {F} } };

  \node[anchor=west, align=left, font=\small] at (4.6,-1.3) {
    \renewcommand{\arraystretch}{1.25}\begin{tabular}{@{}l l@{}}
    \toprule \textbf{Order} & \textbf{Sequence}\\ \midrule
    Preorder (NLR) & A B D E C F\\ Inorder (LNR) & D B E A C F\\ Postorder (LRN) & D E B F C A\\ Level order & A B C D E F\\ \bottomrule
    \end{tabular}};
\end{tikzpicture}
\end{adjustbox}
\caption{Tree traversals. Preorder visits a node before its subtrees, inorder between them and postorder after
them; level order reads the tree row by row.}

\end{figure}


\begin{itemize}
\tightlist
\item
  Unique tree reconstruction needs \textbf{inorder + (preorder or postorder or level-order)}. Preorder + postorder is unique only for full binary trees.
\item
  \textbf{Threaded binary tree}: null right pointers point to inorder successor, null left to predecessor → traversal without stack.
\end{itemize}

\needspace{95pt}

\hypertarget{p2-07--33-binary-search-tree}{%
\subsection{Binary Search Tree}\label{p2-07--33-binary-search-tree}}

\begin{itemize}
\tightlist
\item
  Search/insert/delete: O(h); h = O(log n) average, O(n) worst (skewed).
\item
  Delete node with 2 children: replace with \textbf{inorder successor} (min of right subtree) or predecessor.
\item
  Inorder traversal of BST = \textbf{sorted order}.
\end{itemize}

\needspace{125pt}

\hypertarget{p2-07--34-avl-tree}{%
\subsection{AVL Tree}\label{p2-07--34-avl-tree}}

Balance factor = height(left) − height(right) ∈ \{−1, 0, 1\}.

\begin{figure}[!htbp]
\centering
\begin{adjustbox}{max width=\linewidth}
\begin{tikzpicture}[t/.style={tnode, fill=fslate}, hot/.style={tnode, fill=frose}, e/.style={draw=fgink, line width=0.55pt}]
  % LL
  \node[hot] (a) at (0,0) {30}; \node[t] (b) at (-0.8,-1.1) {20}; \node[t] (c) at (-1.6,-2.2) {10};
  \draw[e] (a) -- (b) -- (c);
  \draw[arr, line width=0.8pt] (0.4,-1.2) -- node[lbl, above] {right rotate} (2.0,-1.2);
  \node[t] (d) at (3.3,-0.4) {20}; \node[t] (e1) at (2.5,-1.5) {10}; \node[t] (f) at (4.1,-1.5) {30};
  \draw[e] (e1) -- (d) -- (f);
  \node[capt] at (1.2,-3.0) {LL case: single right rotation};
  % RR
  \begin{scope}[xshift=7.6cm]
    \node[hot] (a) at (-1.6,0) {10}; \node[t] (b) at (-0.8,-1.1) {20}; \node[t] (c) at (0,-2.2) {30};
    \draw[e] (a) -- (b) -- (c);
    \draw[arr, line width=0.8pt] (0.4,-1.2) -- node[lbl, above] {left rotate} (2.0,-1.2);
    \node[t] (d) at (3.3,-0.4) {20}; \node[t] (e1) at (2.5,-1.5) {10}; \node[t] (f) at (4.1,-1.5) {30};
    \draw[e] (e1) -- (d) -- (f);
    \node[capt] at (1.2,-3.0) {RR case: single left rotation};
  \end{scope}
\end{tikzpicture}
\end{adjustbox}
\caption{AVL single rotations; the unbalanced node is shaded. The double rotations are combinations: LR = left-rotate
the child, then right-rotate the node; RL = right-rotate the child, then left-rotate the node.}
\end{figure}


\begin{itemize}
\tightlist
\item
  Min nodes in AVL of height h: \textbf{N(h) = N(h−1) + N(h−2) + 1}, N(0)=1, N(1)=2 → 1, 2, 4, 7, 12, 20, 33\ldots{}
\item
  Max height ≈ 1.44 log₂ n. All operations O(log n).
\item
  \textbf{Red-Black tree}: root black, no two consecutive reds, same black-height on every path; height ≤ 2 log₂(n+1).
\end{itemize}

\needspace{95pt}

\hypertarget{p2-07--35-b-trees-multi-way-disk-based}{%
\subsection{B-Trees (multi-way, disk-based)}\label{p2-07--35-b-trees-multi-way-disk-based}}

\begin{itemize}
\tightlist
\item
  Order m: each node ≤ m children, ≤ m−1 keys; non-root ≥ ⌈m/2⌉ children; all leaves same level.
\item
  Insertion splits a full node and pushes the median up; tree grows at root.
\item
  \textbf{B+ tree}: all data at leaves, leaves linked. \textbf{B* tree}: nodes at least 2/3 full (redistribution before split).
\item
  Details and order calculations: see \protect\hyperlink{p2-04--8-file-organisation--indexing}{DBMS unit}.
\end{itemize}

\needspace{95pt}

\hypertarget{p2-07--36-heaps}{%
\subsection{Heaps}\label{p2-07--36-heaps}}

\begin{figure}[!htbp]
\centering
\begin{adjustbox}{max width=\linewidth}
\begin{tikzpicture}[level distance=12mm, tn/.style={tnode, fill=fsand},
    level 1/.style={sibling distance=30mm}, level 2/.style={sibling distance=15mm}, edge from parent/.style={draw=fgink, line width=0.55pt}]
  \node[tn] (n0) {90}
    child { node[tn] (n1) {70} child { node[tn] (n3) {30} } child { node[tn] (n4) {60} } }
    child { node[tn] (n2) {80} child { node[tn] (n5) {50} } child[missing] };
  \foreach \i in {0,...,5} \node[font=\scriptsize, text=fwine, inner sep=0pt] at ([xshift=4.6mm, yshift=3.4mm]n\i) {\i};
  \begin{scope}[shift={(4.4,-1.2)}]
    \foreach \v [count=\i from 0] in {90,70,80,30,60,50} {
      \node[draw=fgink, fill=fsand, minimum size=8mm, inner sep=0pt, rectangle, font=\small] at (\i*0.8,0) {\v};
      \node[font=\scriptsize, text=fwine] at (\i*0.8,-0.65) {\i};
    }
    \node[align=left, font=\footnotesize, anchor=north west] at (-0.4,-1.1)
      {parent$(i) = \lfloor (i-1)/2 \rfloor$\\ left$(i) = 2i+1$,\ \ right$(i) = 2i+2$\\ \textit{(0-based indexing)}};
  \end{scope}
\end{tikzpicture}
\end{adjustbox}
\caption{A max-heap and the array that stores it level by level. Every parent is at least as large as its
children, so the maximum sits at index 0.}
\end{figure}


\begin{itemize}
\tightlist
\item
  Insert O(log n); delete-max O(log n); \textbf{build-heap O(n)}; find-max O(1); search O(n).
\item
  In a max-heap, the minimum is among the leaves (⌈n/2⌉ candidates).
\item
  \textbf{Binomial heap} (union O(log n)), \textbf{Fibonacci heap} (decrease-key O(1) amortised → Dijkstra O(E + V log V)).
\end{itemize}

\needspace{191pt}

\hypertarget{p2-07--37-disjoint-sets-unionfind}{%
\subsection{Disjoint Sets (Union--Find)}\label{p2-07--37-disjoint-sets-unionfind}}

Union by rank + path compression → nearly O(1) amortised: \textbf{O(α(n))} (inverse Ackermann). Used in Kruskal.

\needspace{125pt}

\hypertarget{p2-07--4-hashing}{%
\section{Hashing}\label{p2-07--4-hashing}}

Hashing stores a key at an address computed from the key itself, so search, insert and delete take \(O(1)\) time
on average. The design questions are the hash function and what to do when two keys collide.

\begin{formulatable}
Division method ($m$ prime, not close to a power of 2) & h(k) = k \bmod m \\
Multiplication method (Knuth: $A \approx 0.618$) & h(k) = \lfloor m\,(kA \bmod 1) \rfloor \\
Load factor & \alpha = n / m \\
\end{formulatable}


\begingroup\small

\par\medskip\noindent\hfil\begin{tabular}{@{}
  >{\raggedright\arraybackslash}p{(\columnwidth - 4\tabcolsep) * \real{0.2116}}
  >{\raggedright\arraybackslash}p{(\columnwidth - 4\tabcolsep) * \real{0.1868}}
  >{\raggedright\arraybackslash}p{(\columnwidth - 4\tabcolsep) * \real{0.2195}}@{}}
\toprule\noalign{}
\begin{minipage}[b]{\linewidth}\raggedright
\textbf{Collision resolution}
\end{minipage} & \begin{minipage}[b]{\linewidth}\raggedright
\textbf{Probe sequence}
\end{minipage} & \begin{minipage}[b]{\linewidth}\raggedright
\textbf{Issues}
\end{minipage} \\
\midrule\noalign{}
Chaining & Linked list per slot & Expected search 1 + α \\
Linear probing & h(k) + i & \textbf{Primary clustering} \\
Quadratic probing & h(k) + c₁i + c₂i² & \textbf{Secondary clustering} \\
Double hashing & h₁(k) + i·h₂(k) & Best of open addressing \\
\bottomrule\noalign{}
\end{tabular}\hfil\null\par\medskip


\endgroup

Open addressing expected probes: unsuccessful ≤ 1/(1 − α); successful ≤ (1/α) ln(1/(1−α)).

\textbf{Worked (linear probing, m = 10, h(k) = k mod 10)}: insert 12, 18, 13, 2, 3, 23, 5, 15

\begin{figure}[!htbp]
\centering
\begin{adjustbox}{max width=\linewidth}
\begin{tikzpicture}[x=1.05cm]
  \foreach \i/\v/\c in {0/{}/fpaper, 1/{}/fpaper, 2/12/fsage, 3/13/fsage, 4/2/frose, 5/3/frose, 6/23/frose, 7/5/frose, 8/18/fsage, 9/15/frose} {
    \filldraw[fill=\c, draw=fgink, line width=0.5pt] (\i,0) rectangle (\i+1,0.8);
    \node[font=\small] at (\i+0.5,0.4) {\v};
    \node[font=\scriptsize, text=fgink!70] at (\i+0.5,1.05) {\i};
  }
  \node[font=\footnotesize, align=left, anchor=north west] at (0,-0.25) {%
    \renewcommand{\arraystretch}{1.15}\begin{tabular}{@{}l l l l@{}}
    12 $\to$ 2 & 18 $\to$ 8 & 13 $\to$ 3 & \\
    2 $\to$ 2\,\texttimes, 3\,\texttimes, \textbf{4} & 3 $\to$ 3\,\texttimes, 4\,\texttimes, \textbf{5} & 23 $\to$ 3, 4, 5\,\texttimes, \textbf{6} & 5 $\to$ 5, 6\,\texttimes, \textbf{7} \\
    15 $\to$ 5, 6, 7, 8\,\texttimes, \textbf{9} & & & \\
    \end{tabular}};
  \fill[fsage] (7.2,-1.6) rectangle (7.5,-1.35); \node[font=\scriptsize, anchor=west] at (7.55,-1.475) {at home slot};
  \fill[frose] (7.2,-2.0) rectangle (7.5,-1.75); \node[font=\scriptsize, anchor=west] at (7.55,-1.875) {displaced by probing};
\end{tikzpicture}
\end{adjustbox}
\caption{Linear probing, $m = 10$, $h(k) = k \bmod 10$, inserting 12, 18, 13, 2, 3, 23, 5, 15. Keys pile up in
the run 2--9: this is \emph{primary clustering}.}
\end{figure}


\needspace{160pt}

\hypertarget{p2-07--5-sorting}{%
\section{Sorting}\label{p2-07--5-sorting}}

Sorting questions test three things: time complexity in the best, average and worst cases; whether the sort is
stable and in place; and the state of the array after one pass.

\begingroup\small

\par\medskip\noindent\hfil\begin{tabular}{@{}
  >{\raggedright\arraybackslash}p{(\columnwidth - 12\tabcolsep) * \real{0.1357}}
  >{\raggedright\arraybackslash}p{(\columnwidth - 12\tabcolsep) * \real{0.1520}}
  >{\raggedright\arraybackslash}p{(\columnwidth - 12\tabcolsep) * \real{0.1096}}
  >{\raggedright\arraybackslash}p{(\columnwidth - 12\tabcolsep) * \real{0.2133}}
  >{\raggedright\arraybackslash}p{(\columnwidth - 12\tabcolsep) * \real{0.0879}}
  >{\raggedright\arraybackslash}p{(\columnwidth - 12\tabcolsep) * \real{0.0954}}
  >{\raggedright\arraybackslash}p{(\columnwidth - 12\tabcolsep) * \real{0.1215}}@{}}
\toprule\noalign{}
\begin{minipage}[b]{\linewidth}\raggedright
\textbf{Algorithm}
\end{minipage} & \begin{minipage}[b]{\linewidth}\raggedright
\textbf{Best}
\end{minipage} & \begin{minipage}[b]{\linewidth}\raggedright
\textbf{Average}
\end{minipage} & \begin{minipage}[b]{\linewidth}\raggedright
\textbf{Worst}
\end{minipage} & \begin{minipage}[b]{\linewidth}\raggedright
\textbf{Space}
\end{minipage} & \begin{minipage}[b]{\linewidth}\raggedright
\textbf{Stable}
\end{minipage} & \begin{minipage}[b]{\linewidth}\raggedright
\textbf{In-place}
\end{minipage} \\
\midrule\noalign{}
Bubble & n (with flag) & n² & n² & 1 & ✔ & ✔ \\
Selection & n² & n² & n² & 1 & ✘ & ✔ \\
Insertion & n & n² & n² & 1 & ✔ & ✔ \\
Merge & n log n & n log n & n log n & n & ✔ & ✘ \\
\textbf{Quick} & n log n & n log n & \textbf{n²} & log n & ✘ & ✔ \\
Heap & n log n & n log n & n log n & 1 & ✘ & ✔ \\
Shell & n log n & \textasciitilde n\^{}1.3 & n² (gap-dependent) & 1 & ✘ & ✔ \\
Counting & n + k & n + k & n + k & k & ✔ & ✘ \\
Radix & d(n + k) & d(n + k) & d(n + k) & n + k & ✔ & ✘ \\
Bucket & n + k & n + k & n² & n & ✔ & ✘ \\
\bottomrule\noalign{}
\end{tabular}\hfil\null\par\medskip


\endgroup

\begin{itemize}
\tightlist
\item
  Selection sort makes the \textbf{minimum number of swaps} (n − 1).
\item
  Insertion sort is best for \textbf{nearly sorted} data; number of comparisons ≈ number of inversions.
\item
  Quick sort worst case on sorted input with first/last element as pivot. Randomised pivot → expected O(n log n).
\item
  \textbf{Comparison-sort lower bound: Ω(n log n)} (decision tree with n! leaves has height ≥ log₂ n!).
\item
  Merging k sorted lists of total n: O(n log k) with min-heap.
\end{itemize}

\begin{figure}[!htbp]
\centering
\begin{adjustbox}{max width=\linewidth}
\begin{tikzpicture}[x=0.85cm]
  \foreach \v/\c [count=\i from 0] in {7/fpaper, 2/fpaper, 1/fpaper, 6/fpaper, 8/fpaper, 5/fpaper, 3/fpaper, 4/fochre} {
    \filldraw[fill=\c, draw=fgink, line width=0.5pt] (\i,0) rectangle (\i+1,0.8); \node[font=\small] at (\i+0.5,0.4) {\v}; }
  \node[capt, anchor=east] at (-0.2,0.4) {before};
  \node[lbl, anchor=west] at (8.2,0.4) {pivot $= 4$ (last element)};
  \foreach \v/\c [count=\i from 0] in {2/fsage, 1/fsage, 3/fsage, 4/fochre, 8/frose, 5/frose, 7/frose, 6/frose} {
    \filldraw[fill=\c, draw=fgink, line width=0.5pt] (\i,-1.8) rectangle (\i+1,-1.0); \node[font=\small] at (\i+0.5,-1.4) {\v}; }
  \node[capt, anchor=east] at (-0.2,-1.4) {after};
  \draw[brace] (3,-1.95) -- node[lbl, below=4pt] {$\le$ pivot} (0,-1.95);
  \draw[brace] (8,-1.95) -- node[lbl, below=4pt] {$>$ pivot} (4,-1.95);
  \draw[arr, fgink!60] (7.5,-0.05) to[bend left=25] (3.5,-0.95);
  \node[lbl, anchor=west] at (8.2,-1.4) {pivot in its final place};
\end{tikzpicture}
\end{adjustbox}
\caption{Lomuto partition. One pass places the pivot in its final position, with smaller keys to its left and
larger keys to its right; quicksort then recurses on the two sides.}
\end{figure}


\needspace{95pt}

\hypertarget{p2-07--6-searching}{%
\section{Searching}\label{p2-07--6-searching}}

\begin{itemize}
\tightlist
\item
  Linear O(n); \textbf{Binary} O(log n) on sorted array --- max comparisons ⌊log₂ n⌋ + 1.
\item
  Interpolation search O(log log n) average on uniform data.
\item
  Selection (k-th smallest): Quickselect O(n) average; \textbf{median of medians O(n) worst}.
\item
  Min \& max together: \textbf{3n/2 − 2} comparisons. Second largest: n + ⌈log₂ n⌉ − 2.
\end{itemize}

\needspace{125pt}

\hypertarget{p2-07--7-algorithm-design-techniques}{%
\section{Algorithm Design Techniques}\label{p2-07--7-algorithm-design-techniques}}

Most algorithms in the syllabus follow one of five strategies. Recognising the strategy tells you the recurrence
and therefore the running time.

\begin{figure}[!htbp]
\centering
\begin{adjustbox}{max width=\linewidth}
\begin{tikzpicture}[card/.style={ink, rectangle, rounded corners=2pt, fill=fpaper, text width=46mm, align=left, inner sep=5pt, font=\footnotesize, anchor=north,
    blur shadow={shadow blur steps=6, shadow xshift=0.8pt, shadow yshift=-1pt, shadow opacity=22}},
    hd/.style={font=\small\bfseries, anchor=south, inner sep=2pt}]
  \foreach \x/\y/\t/\c/\b in {0/0/Divide \& conquer/fslate/{merge sort, quicksort, binary search, Strassen, closest pair, Karatsuba},
      5.4/0/Greedy/fsage/{Kruskal, Prim, Dijkstra, Huffman, fractional knapsack, activity selection, job sequencing},
      10.8/0/Dynamic programming/fsand/{0/1 knapsack, LCS, matrix chain, Floyd--Warshall, Bellman--Ford, OBST, TSP (Held--Karp)},
      2.7/-3.3/Backtracking/frose/{N-queens, sum of subsets, graph colouring, Hamiltonian cycle, Sudoku},
      8.1/-3.3/Branch \& bound/fgrey/{0/1 knapsack, TSP, job assignment; prunes with bounds, BFS / least-cost search}} {
    \node[card] (k) at (\x,\y) {\b};
    \fill[\c] (k.north west) rectangle ([yshift=6.5mm]k.north east);
    \draw[fgink, line width=0.55pt, rounded corners=2pt] ([yshift=6.5mm]k.north west) rectangle (k.south east);
    \draw[fgink, line width=0.4pt] (k.north west) -- (k.north east);
    \node[hd] at (k.north) {\t};
  }
\end{tikzpicture}
\end{adjustbox}
\caption{The five algorithm-design paradigms and the standard problems the examination attaches to each.}
\end{figure}


\begingroup\small

\par\medskip\noindent\hfil\begin{tabular}{@{}
  >{\raggedright\arraybackslash}p{(\columnwidth - 2\tabcolsep) * \real{0.3983}}
  >{\raggedright\arraybackslash}p{(\columnwidth - 2\tabcolsep) * \real{0.4139}}@{}}
\toprule\noalign{}
\begin{minipage}[b]{\linewidth}\raggedright
\textbf{Greedy}
\end{minipage} & \begin{minipage}[b]{\linewidth}\raggedright
\textbf{Dynamic programming}
\end{minipage} \\
\midrule\noalign{}
Locally optimal choice, never reconsidered & Solves overlapping subproblems, stores results \\
\textbf{Greedy-choice property} + optimal substructure & \textbf{Optimal substructure + overlapping subproblems} \\
Faster, not always optimal & Always optimal (if formulated correctly) \\
Fractional knapsack ✔ ; 0/1 knapsack ✘ & 0/1 knapsack ✔ O(nW) --- pseudo-polynomial \\
\bottomrule\noalign{}
\end{tabular}\hfil\null\par\medskip


\endgroup

\needspace{125pt}

\hypertarget{p2-07--71-key-dp-problems}{%
\subsection{Key DP Problems}\label{p2-07--71-key-dp-problems}}

\textbf{LCS} of X = ABCBDAB, Y = BDCABA → length \textbf{4} (e.g. BCBA). Time O(mn).

\[
\text{LCS}[i][j] =
\begin{cases}
\text{LCS}[i-1][j-1] + 1 & \text{if } x_i = y_j\\[2pt]
\max\bigl(\text{LCS}[i-1][j],\ \text{LCS}[i][j-1]\bigr) & \text{otherwise}
\end{cases}
\]


\textbf{Matrix chain multiplication}: dims 10×30, 30×5, 5×60.

\begin{formulas}
\begin{align*}
(AB)C &= 10\cdot30\cdot5 + 10\cdot5\cdot60 = 1500 + 3000 = \mathbf{4500} \quad\text{(optimal)}\\
A(BC) &= 30\cdot5\cdot60 + 10\cdot30\cdot60 = 9000 + 18000 = 27000\\
m[i][j] &= \min_{i \le k < j}\bigl(m[i][k] + m[k+1][j] + p_{i-1}\,p_k\,p_j\bigr) \qquad O(n^3)
\end{align*}
\end{formulas}
\noindent The number of ways to parenthesise a product of $n$ matrices is the Catalan number $C_{n-1}$.


\textbf{0/1 Knapsack}: K{[}i{]}{[}w{]} = max(K{[}i−1{]}{[}w{]}, vᵢ + K{[}i−1{]}{[}w − wᵢ{]}).

\needspace{95pt}

\hypertarget{p2-07--72-greedy-problems}{%
\subsection{Greedy Problems}\label{p2-07--72-greedy-problems}}

\begin{itemize}
\tightlist
\item
  \textbf{Huffman coding}: repeatedly merge two least-frequent nodes. O(n log n).
  Frequencies a:5, b:9, c:12, d:13, e:16, f:45 → codes f:0, c:100, d:101, a:1100, b:1101, e:111. Average bits = 2.24.
\item
  \textbf{Activity selection}: sort by finish time, pick compatible.
\item
  \textbf{Job sequencing with deadlines}: sort by profit, place in latest free slot.
\item
  \textbf{Fractional knapsack}: sort by value/weight.
\item
  \textbf{Optimal merge pattern}: like Huffman.
\end{itemize}

\needspace{95pt}

\hypertarget{p2-07--73-backtracking}{%
\subsection{Backtracking}\label{p2-07--73-backtracking}}

\begin{itemize}
\tightlist
\item
  N-Queens: 4-Queens has \textbf{2} solutions; 8-Queens has \textbf{92}.
\item
  State-space tree; prune when constraint violated (bounding function).
\item
  Graph m-colouring, sum of subsets, Hamiltonian cycles.
\end{itemize}

\needspace{95pt}

\hypertarget{p2-07--74-branch--bound}{%
\subsection{Branch \& Bound}\label{p2-07--74-branch--bound}}

\begin{itemize}
\tightlist
\item
  Uses BFS / \textbf{least-cost (LC) search} with bounds; FIFO B\&B, LIFO B\&B, LC B\&B. Used for optimisation (TSP, 0/1 knapsack, assignment).
\end{itemize}

\needspace{227pt}

\hypertarget{p2-07--8-graph-algorithms}{%
\section{Graph Algorithms}\label{p2-07--8-graph-algorithms}}

A graph \(G = (V, E)\) models any network of relationships. The standard algorithms traverse it, find shortest
paths, build minimum spanning trees and compute maximum flows.

\needspace{161pt}

\hypertarget{p2-07--81-representations}{%
\subsection{Representations}\label{p2-07--81-representations}}

Adjacency matrix: O(V²) space, O(1) edge check. Adjacency list: O(V + E).

\needspace{95pt}

\hypertarget{p2-07--82-bfs--dfs}{%
\subsection{BFS \& DFS}\label{p2-07--82-bfs--dfs}}

\begin{figure}[!htbp]
\centering
\begin{adjustbox}{max width=\linewidth}
\begin{tikzpicture}[v/.style={vertex, fill=fslate, minimum size=8mm}]
  \node[v] (A) at (0,1.6) {A}; \node[v] (B) at (2,1.6) {B}; \node[v] (E) at (4,1.6) {E};
  \node[v] (C) at (0,0) {C}; \node[v] (D) at (2,0) {D};
  \draw[thinline] (A) -- (B) -- (E) (A) -- (C) -- (D) -- (B);
  \node[font=\footnotesize, anchor=west, align=left] at (5.2,0.8) {%
    \renewcommand{\arraystretch}{1.3}\begin{tabular}{@{}l l l@{}}
    \toprule & \textbf{Order from A} & \textbf{Uses} \\ \midrule
    BFS (queue) & A B C E D & shortest paths in unweighted graphs \\
    DFS (stack) & A B E D C & topological sort, SCCs, cycle detection \\
    \bottomrule \end{tabular}};
\end{tikzpicture}
\end{adjustbox}
\caption{BFS explores level by level; DFS goes as deep as possible before backtracking. Both take $O(V+E)$ with
adjacency lists and $O(V^2)$ with a matrix. (The exact order depends on the order in which neighbours are listed.)}
\end{figure}


\begin{itemize}
\tightlist
\item
  DFS edge types: tree, back (→ cycle), forward, cross. Undirected DFS has only tree and back edges.
\item
  \textbf{Topological sort} (DAG only): DFS finishing order reversed, or Kahn\textquotesingle s algorithm (in-degree 0).
\item
  \textbf{Strongly connected components}: Kosaraju (2 DFS), Tarjan (1 DFS). O(V + E).
\item
  Articulation points \& bridges: DFS with low values.
\end{itemize}

\needspace{130pt}

\hypertarget{p2-07--83-shortest-paths}{%
\subsection{Shortest Paths}\label{p2-07--83-shortest-paths}}

\begingroup\small

\par\medskip\noindent\hfil\begin{tabular}{@{}
  >{\raggedright\arraybackslash}p{(\columnwidth - 6\tabcolsep) * \real{0.1863}}
  >{\raggedright\arraybackslash}p{(\columnwidth - 6\tabcolsep) * \real{0.2577}}
  >{\raggedright\arraybackslash}p{(\columnwidth - 6\tabcolsep) * \real{0.2732}}
  >{\raggedright\arraybackslash}p{(\columnwidth - 6\tabcolsep) * \real{0.2827}}@{}}
\toprule\noalign{}
\begin{minipage}[b]{\linewidth}\raggedright
\textbf{Algorithm}
\end{minipage} & \begin{minipage}[b]{\linewidth}\raggedright
\textbf{Type}
\end{minipage} & \begin{minipage}[b]{\linewidth}\raggedright
\textbf{Negative edges?}
\end{minipage} & \begin{minipage}[b]{\linewidth}\raggedright
\textbf{Complexity}
\end{minipage} \\
\midrule\noalign{}
BFS & Single-source, unweighted & --- & O(V + E) \\
\textbf{Dijkstra} & Single-source & \textbf{No} & O((V + E) log V) binary heap; O(V²) array; O(E + V log V) Fibonacci \\
\textbf{Bellman--Ford} & Single-source & Yes; detects negative cycles & O(VE) \\
DAG shortest path & Single-source & Yes & O(V + E) \\
\textbf{Floyd--Warshall} & All-pairs (DP) & Yes (no neg. cycles) & O(V³) \\
Johnson & All-pairs, sparse & Yes & O(V² log V + VE) \\
\bottomrule\noalign{}
\end{tabular}\hfil\null\par\medskip


\endgroup

\textbf{Dijkstra worked} (directed edges: A→B 4, A→C 1, C→B 2, C→D 5, B→D 1; source A)

\begin{figure}[!htbp]
\centering
\begin{adjustbox}{max width=\linewidth}
\begin{tikzpicture}[v/.style={vertex, fill=fslate, minimum size=8mm}, w/.style={edgelbl}]
  \node[v, fill=fochre] (A) at (0,2) {A}; \node[v] (B) at (3,2) {B};
  \node[v] (C) at (0,0) {C}; \node[v] (D) at (3,0) {D};
  \draw[arr] (A) -- node[w] {4} (B);
  \draw[arr, fwine, line width=1pt] (A) -- node[w] {1} (C);
  \draw[arr, fwine, line width=1pt] (C) -- node[w] {2} (B);
  \draw[arr] (C) -- node[w] {5} (D);
  \draw[arr, fwine, line width=1pt] (B) -- node[w] {1} (D);
  \node[smalllbl] at (-0.2,2.65) {source};
  \node[font=\footnotesize, anchor=west] at (4.4,1.0) {%
    \renewcommand{\arraystretch}{1.25}\begin{tabular}{@{}c l c l@{}}
    \toprule \textbf{Step} & \textbf{Visited} & \textbf{dist(A, B, C, D)} & \textbf{Update}\\ \midrule
    0 & --- & 0, $\infty$, $\infty$, $\infty$ & \\
    1 & A & 0, 4, 1, $\infty$ & \\
    2 & A, C & 0, \textbf{3}, 1, 6 & B via C $= 1+2$; D via C $= 1+5$ \\
    3 & A, C, B & 0, 3, 1, \textbf{4} & D via B $= 3+1$ \\
    4 & A, C, B, D & 0, 3, 1, 4 & final \\
    \bottomrule \end{tabular}};
\end{tikzpicture}
\end{adjustbox}
\caption{Dijkstra's algorithm from A. Each step fixes the closest unvisited vertex and relaxes its out-edges; the
shortest-path tree is drawn in red.}
\end{figure}


\needspace{130pt}

\hypertarget{p2-07--84-minimum-spanning-tree}{%
\subsection{Minimum Spanning Tree}\label{p2-07--84-minimum-spanning-tree}}

\begingroup\small

\par\medskip\noindent\hfil\begin{tabular}{@{}
  >{\raggedright\arraybackslash}p{(\columnwidth - 2\tabcolsep) * \real{0.5200}}
  >{\raggedright\arraybackslash}p{(\columnwidth - 2\tabcolsep) * \real{0.4800}}@{}}
\toprule\noalign{}
\begin{minipage}[b]{\linewidth}\raggedright
\textbf{Kruskal}
\end{minipage} & \begin{minipage}[b]{\linewidth}\raggedright
\textbf{Prim}
\end{minipage} \\
\midrule\noalign{}
Sort edges, add smallest that doesn\textquotesingle t form a cycle (union-find) & Grow tree from a vertex, add cheapest edge leaving tree \\
O(E log E) & O(E log V) heap; O(V²) matrix \\
Better for sparse graphs & Better for dense graphs \\
Forest during execution & Always a single tree \\
\bottomrule\noalign{}
\end{tabular}\hfil\null\par\medskip


\endgroup

\begin{itemize}
\tightlist
\item
  MST is unique if all edge weights are distinct.
\item
  \textbf{Cut property}: lightest edge across any cut belongs to some MST. \textbf{Cycle property}: heaviest edge in a cycle is not in any MST.
\item
  Number of spanning trees of Kₙ: nⁿ⁻².
\end{itemize}

\needspace{95pt}

\hypertarget{p2-07--85-maximum-flow}{%
\subsection{Maximum Flow}\label{p2-07--85-maximum-flow}}

\begin{itemize}
\tightlist
\item
  \textbf{Ford--Fulkerson} (augmenting paths): O(E · f*). \textbf{Edmonds--Karp} (BFS augmenting paths): O(VE²).
\item
  \textbf{Max-flow min-cut theorem}: value of max flow = capacity of min s--t cut.
\item
  Bipartite matching reduces to max flow.
\end{itemize}

\needspace{95pt}

\hypertarget{p2-07--9-lower-bound-theory}{%
\section{Lower Bound Theory}\label{p2-07--9-lower-bound-theory}}

\begin{itemize}
\tightlist
\item
  \textbf{Comparison tree (decision tree)}: sorting needs ≥ ⌈log₂ n!⌉ = Ω(n log n) comparisons; searching sorted array Ω(log n).
\item
  Merging two sorted lists of size n: 2n − 1 comparisons (worst). Finding max: n − 1 (tournament argument).
\item
  \textbf{Reductions}: if A reduces to B and A has lower bound L, B has lower bound L (e.g. sorting ≤ convex hull → convex hull Ω(n log n)).
\end{itemize}

\needspace{125pt}

\hypertarget{p2-07--10-complexity-classes}{%
\section{Complexity Classes}\label{p2-07--10-complexity-classes}}

Complexity classes group problems, not algorithms, by the resources the best possible algorithm needs. The
open question P = NP asks whether every problem whose solution can be checked quickly can also be solved quickly.

\begin{figure}[!htbp]
\centering
\begin{adjustbox}{max width=\linewidth}
\begin{tikzpicture}
  \begin{scope}
    \clip (0,0) ellipse (3.6 and 2.2);
    \fill[fsand!70] (3.2,0.0) ellipse (2.5 and 2.6);
  \end{scope}
  \draw[fgink, line width=0.6pt, fill=fslate!35, fill opacity=0.0] (0,0) ellipse (3.6 and 2.2);
  \filldraw[fill=fsage, draw=fgink, line width=0.6pt] (-1.3,0) ellipse (1.5 and 1.1);
  \draw[fgink, line width=0.6pt, dash pattern=on 4pt off 2pt] (3.2,0.0) ellipse (2.5 and 2.6);
  \node[font=\small, align=center] at (-1.3,0) {\textbf{P}\\\footnotesize solvable in\\\footnotesize polynomial time};
  \node[font=\small, align=center] at (1.75,0) {\textbf{NP-}\\\textbf{complete}};
  \node[font=\small, align=center, anchor=south] at (-1.2,2.25) {\textbf{NP}: verifiable in polynomial time};
  \node[font=\small, align=left, anchor=west] at (5.8,1.2) {\textbf{NP-hard}\\\footnotesize at least as hard\\\footnotesize as every problem in NP};
\end{tikzpicture}
\end{adjustbox}
\caption{Complexity classes, assuming P $\ne$ NP. NP-complete = NP $\cap$ NP-hard. If any NP-complete problem were
solvable in polynomial time, then P = NP.}
\end{figure}


\begin{itemize}
\tightlist
\item
  \textbf{P ⊆ NP} (P = NP? is open). NP-Complete = NP ∩ NP-Hard.
\item
  To prove X is NP-complete: (1) X ∈ NP; (2) reduce a known NP-complete problem \textbf{to X} (Y ≤ₚ X).
\item
  \textbf{Cook--Levin theorem}: \textbf{SAT} is the first NP-complete problem.
\item
  NP-complete: SAT, 3-SAT, CLIQUE, Vertex Cover, Independent Set, Hamiltonian cycle, TSP (decision), Subset sum, 0/1 Knapsack (decision), Graph colouring (k ≥ 3), Set cover, Partition.
\item
  In P: 2-SAT, shortest path, MST, 2-colouring (bipartite test), Euler circuit, sorting, matching, linear programming, primality (AKS 2002).
\item
  \textbf{NP-Hard but not in NP}: Halting problem; optimisation TSP.
\item
  Reductions chain: SAT → 3-SAT → CLIQUE → Vertex Cover → Hamiltonian cycle → TSP.
\end{itemize}

\needspace{130pt}

\hypertarget{p2-07--11-selected-topics}{%
\section{Selected Topics}\label{p2-07--11-selected-topics}}

\begingroup\small

\par\medskip\noindent\hfil\begin{tabular}{@{}
  >{\raggedright\arraybackslash}p{(\columnwidth - 2\tabcolsep) * \real{0.2389}}
  >{\raggedright\arraybackslash}p{(\columnwidth - 2\tabcolsep) * \real{0.7611}}@{}}
\toprule\noalign{}
\begin{minipage}[b]{\linewidth}\raggedright
\textbf{Topic}
\end{minipage} & \begin{minipage}[b]{\linewidth}\raggedright
\textbf{Key result}
\end{minipage} \\
\midrule\noalign{}
GCD (Euclid) & gcd(a, b) = gcd(b, a mod b); O(log min(a,b)); extended Euclid gives x, y with ax + by = gcd \\
Modular exponentiation & Repeated squaring O(log e) \\
RSA & n = pq, φ = (p−1)(q−1), ed ≡ 1 mod φ \\
Primality & Miller--Rabin (randomised), AKS (deterministic polynomial) \\
Polynomial multiplication & Naive O(n²); \textbf{FFT O(n log n)} (evaluate at roots of unity, pointwise multiply, inverse FFT) \\
Naive string match & O((n − m + 1)m) \\
\textbf{Rabin--Karp} & Rolling hash; average O(n + m), worst O(nm) \\
\textbf{KMP} & Prefix (failure) function; \textbf{O(n + m)} \\
Boyer--Moore & Bad-character \& good-suffix; sub-linear in practice \\
Finite automaton matcher & O(n) matching after O(m· \\
\bottomrule\noalign{}
\end{tabular}\hfil\null\par\medskip


\endgroup

KMP prefix function for pattern \textbf{"ababaca"}: π = {[}0, 0, 1, 2, 3, 0, 1{]}.

\needspace{95pt}

\hypertarget{p2-07--12-advanced-algorithms}{%
\section{Advanced Algorithms}\label{p2-07--12-advanced-algorithms}}

\begin{itemize}
\tightlist
\item
  \textbf{Parallel algorithms} (PRAM models: EREW, CREW, CRCW). Parallel prefix sum O(log n) with n processors. Odd-even transposition sort O(n) with n processors; bitonic sort O(log² n). Parallel merge O(log n). Brent\textquotesingle s theorem.
\item
  \textbf{Approximation algorithms}: Vertex cover --- 2-approximation (pick both ends of uncovered edge); Metric TSP --- 2-approx (MST doubling), \textbf{1.5-approx (Christofides)}; Set cover --- ln n approx (greedy). Approximation ratio ρ(n).
\item
  \textbf{Randomised}: \textbf{Las Vegas} (always correct, random time --- randomised quicksort) vs \textbf{Monte Carlo} (fixed time, may be wrong --- Miller--Rabin, Karger\textquotesingle s min-cut).
\end{itemize}

\needspace{131pt}

\hypertarget{p2-07--13-deeper-dive--traces-you-must-be-able-to-do-by-hand}{%
\section{Deeper Dive --- Traces You Must Be Able to Do by Hand}\label{p2-07--13-deeper-dive--traces-you-must-be-able-to-do-by-hand}}

\needspace{95pt}

\hypertarget{p2-07--131-recursion-tree-for-tn--2tn2--n}{%
\subsection{Recursion Tree for T(n) = 2T(n/2) + n}\label{p2-07--131-recursion-tree-for-tn--2tn2--n}}

\begin{figure}[!htbp]
\centering
\begin{adjustbox}{max width=\linewidth}
\begin{tikzpicture}[r/.style={ink, rectangle, rounded corners=1.2pt, fill=fslate, minimum width=11mm, minimum height=6.5mm, font=\footnotesize, inner sep=1pt}]
  \node[r] (a) at (0,0) {$n$};
  \node[r] (b1) at (-2.4,-1.2) {$n/2$}; \node[r] (b2) at (2.4,-1.2) {$n/2$};
  \foreach \x [count=\i] in {-3.6,-1.2,1.2,3.6} \node[r] (c\i) at (\x,-2.4) {$n/4$};
  \draw[thinline] (a) -- (b1) (a) -- (b2) (b1) -- (c1) (b1) -- (c2) (b2) -- (c3) (b2) -- (c4);
  \foreach \i in {1,...,4} { \draw[thinline, dash pattern=on 1.5pt off 1.5pt] (c\i) -- ++(-0.45,-0.75); \draw[thinline, dash pattern=on 1.5pt off 1.5pt] (c\i) -- ++(0.45,-0.75); }
  \foreach \x in {-4.2,-3.6,...,4.2} \node[r, fill=fsand, minimum width=4mm] at (\x,-4.2) {\scriptsize 1};
  \node at (0,-3.45) {$\vdots$};
  \foreach \y/\t in {0/{$n$}, -1.2/{$2\cdot n/2 = n$}, -2.4/{$4\cdot n/4 = n$}, -4.2/{$n \cdot 1 = n$}}
    \node[font=\footnotesize, anchor=west] at (5.3,\y) {\t};
  \node[capt, anchor=west] at (5.3,0.65) {cost per level};
  \draw[brace] (-4.8,-4.2) -- node[lbl, left=5pt] {$\log_2 n + 1$\\levels} (-4.8,0);
\end{tikzpicture}
\end{adjustbox}
\caption{Recursion tree for $T(n) = 2T(n/2) + n$. Every level costs $n$ and there are $\log_2 n + 1$ levels, so
$T(n) = n(\log_2 n + 1) = \Theta(n\log n)$.}
\end{figure}


For T(n) = T(n − 1) + n: unrolling gives n + (n − 1) + \ldots{} + 1 = n(n + 1)/2 = \textbf{Θ(n²)}.

\needspace{95pt}

\hypertarget{p2-07--132-avl-insertion-trace-10-20-30-40-50-25}{%
\subsection{AVL Insertion Trace: 10, 20, 30, 40, 50, 25}\label{p2-07--132-avl-insertion-trace-10-20-30-40-50-25}}

\begin{figure}[!htbp]
\centering
\begin{adjustbox}{max width=\linewidth}
\begin{tikzpicture}[t/.style={tnode, fill=fslate, minimum size=7mm, font=\footnotesize}, hot/.style={t, fill=frose}, new/.style={t, fill=fsand},
    e/.style={draw=fgink, line width=0.5pt}, st/.style={capt, align=center}]
  % (a) after 10,20,30
  \node[t] (a1) at (0,0) {20}; \node[t] (a2) at (-0.7,-1) {10}; \node[t] (a3) at (0.7,-1) {30};
  \draw[e] (a2) -- (a1) -- (a3);
  \node[st] at (0,-2.1) {(a) 10, 20, 30\\RR at 10: rotate left};
  % (b) after 40,50
  \begin{scope}[xshift=3.9cm]
    \node[t] (b1) at (0,0) {20}; \node[t] (b2) at (-0.7,-1) {10}; \node[t] (b3) at (0.7,-1) {40};
    \node[t] (b4) at (0.1,-2) {30}; \node[t] (b5) at (1.3,-2) {50};
    \draw[e] (b2) -- (b1) -- (b3) (b4) -- (b3) -- (b5);
    \node[st] at (0.3,-3.1) {(b) 40, 50\\RR at 30: rotate left};
  \end{scope}
  % (c) insert 25: imbalance at 20
  \begin{scope}[xshift=8.2cm]
    \node[hot] (c1) at (0,0) {20}; \node[t] (c2) at (-0.7,-1) {10}; \node[t] (c3) at (0.7,-1) {40};
    \node[t] (c4) at (0.1,-2) {30}; \node[t] (c5) at (1.3,-2) {50}; \node[new] (c6) at (-0.5,-3) {25};
    \draw[e] (c2) -- (c1) -- (c3) (c4) -- (c3) -- (c5) (c6) -- (c4);
    \node[st] at (0.3,-4.0) {(c) insert 25:\\20 has balance $-2$ (RL)};
  \end{scope}
  % (d) right-rotate 40
  \begin{scope}[xshift=12.5cm]
    \node[hot] (d1) at (0,0) {20}; \node[t] (d2) at (-0.7,-1) {10}; \node[t] (d3) at (0.7,-1) {30};
    \node[new] (d4) at (0.1,-2) {25}; \node[t] (d5) at (1.3,-2) {40}; \node[t] (d6) at (1.9,-3) {50};
    \draw[e] (d2) -- (d1) -- (d3) (d4) -- (d3) -- (d5) -- (d6);
    \node[st] at (0.5,-4.0) {(d) step 1:\\right-rotate 40};
  \end{scope}
  % (e) left-rotate 20
  \begin{scope}[xshift=17.2cm]
    \node[t] (e1) at (0,0) {30}; \node[t] (e2) at (-1.0,-1) {20}; \node[t] (e3) at (1.0,-1) {40};
    \node[t] (e4) at (-1.6,-2) {10}; \node[new] (e5) at (-0.4,-2) {25}; \node[t] (e6) at (1.6,-2) {50};
    \draw[e] (e2) -- (e1) -- (e3) (e4) -- (e2) -- (e5) (e3) -- (e6);
    \node[st] at (0,-4.0) {(e) step 2:\\left-rotate 20};
  \end{scope}
\end{tikzpicture}
\end{adjustbox}
\caption{AVL insertion of 10, 20, 30, 40, 50, 25. The unbalanced node is shaded red, the newly inserted key yellow.
The final RL case needs a double rotation.}
\end{figure}


\needspace{130pt}

\hypertarget{p2-07--133-heap-sort-trace-4-10-3-5-1}{%
\subsection{Heap Sort Trace: {[}4, 10, 3, 5, 1{]}}\label{p2-07--133-heap-sort-trace-4-10-3-5-1}}

\begingroup\small

\par\medskip\noindent\hfil\begin{tabular}{@{}
  >{\raggedright\arraybackslash}p{(\columnwidth - 4\tabcolsep) * \real{0.2121}}
  >{\raggedright\arraybackslash}p{(\columnwidth - 4\tabcolsep) * \real{0.5207}}
  >{\raggedright\arraybackslash}p{(\columnwidth - 4\tabcolsep) * \real{0.2149}}@{}}
\toprule\noalign{}
\begin{minipage}[b]{\linewidth}\raggedright
\textbf{Step}
\end{minipage} & \begin{minipage}[b]{\linewidth}\raggedright
\textbf{Array}
\end{minipage} & \begin{minipage}[b]{\linewidth}\raggedright
\textbf{Note}
\end{minipage} \\
\midrule\noalign{}
Build heap & {[}4, 10, 3, 5, 1{]} → {[}10, 4, 3, 5, 1{]} → {[}10, 5, 3, 4, 1{]} & heapify index 1, then 0 \\
Swap 10 ↔ 1, heapify & {[}5, 4, 3, 1 \textbar{} 10{]} & \\
Swap 5 ↔ 1, heapify & {[}4, 1, 3 \textbar{} 5, 10{]} & \\
Swap 4 ↔ 3, heapify & {[}3, 1 \textbar{} 4, 5, 10{]} & \\
Swap 3 ↔ 1 & {[}1 \textbar{} 3, 4, 5, 10{]} & sorted \\
\bottomrule\noalign{}
\end{tabular}\hfil\null\par\medskip


\endgroup

\needspace{95pt}

\hypertarget{p2-07--134-kruskal-vs-prim-on-the-same-graph}{%
\subsection{Kruskal vs Prim on the Same Graph}\label{p2-07--134-kruskal-vs-prim-on-the-same-graph}}

\begin{figure}[!htbp]
\centering
\begin{adjustbox}{max width=\linewidth}
\begin{tikzpicture}[v/.style={vertex, fill=fslate, minimum size=8mm}, w/.style={edgelbl}, mst/.style={draw=fwine, line width=1.4pt}]
  \node[v] (A) at (0,1.8) {A}; \node[v] (B) at (3,1.8) {B}; \node[v] (C) at (0,0) {C};
  \node[v] (D) at (3,0) {D}; \node[v] (E) at (5.6,0.9) {E};
  \draw[thinline] (A) -- node[w] {4} (B);
  \draw[mst] (A) -- node[w] {1} (C);
  \draw[mst] (B) -- node[w] {2} (C);
  \draw[mst] (B) -- node[w] {5} (D);
  \draw[thinline] (C) -- node[w] {8} (D);
  \draw[mst] (D) -- node[w] {3} (E);
  \draw[thinline] (C) to[bend right=40] node[w] {9} (E);
\end{tikzpicture}
\end{adjustbox}
\caption{The minimum spanning tree (thick red) has weight $1 + 2 + 3 + 5 = 11$. Kruskal and Prim add the same
edges, in different orders.}
\end{figure}


\begingroup\small

\par\medskip\noindent\hfil\begin{tabular}{@{}
  >{\raggedright\arraybackslash}p{(\columnwidth - 6\tabcolsep) * \real{0.2452}}
  >{\raggedright\arraybackslash}p{(\columnwidth - 6\tabcolsep) * \real{0.1256}}
  >{\raggedright\arraybackslash}p{(\columnwidth - 6\tabcolsep) * \real{0.1732}}
  >{\raggedright\arraybackslash}p{(\columnwidth - 6\tabcolsep) * \real{0.1285}}@{}}
\toprule\noalign{}
\begin{minipage}[b]{\linewidth}\raggedright
\textbf{Kruskal (sorted edges)}
\end{minipage} & \begin{minipage}[b]{\linewidth}\raggedright
\textbf{Decision}
\end{minipage} & \begin{minipage}[b]{\linewidth}\raggedright
\textbf{Prim (start A)}
\end{minipage} & \begin{minipage}[b]{\linewidth}\raggedright
\textbf{Edge added}
\end{minipage} \\
\midrule\noalign{}
A--C 1 & take & \{A\} & A--C 1 \\
B--C 2 & take & \{A, C\} & B--C 2 \\
D--E 3 & take & \{A, B, C\} & B--D 5 \\
A--B 4 & skip (cycle) & \{A, B, C, D\} & D--E 3 \\
B--D 5 & take → done & & \\
\bottomrule\noalign{}
\end{tabular}\hfil\null\par\medskip


\endgroup

Both give MST weight \textbf{1 + 2 + 3 + 5 = 11}.

\needspace{125pt}

\hypertarget{p2-07--135-floydwarshall-trace}{%
\subsection{Floyd--Warshall Trace}\label{p2-07--135-floydwarshall-trace}}

Directed graph: 1→2 (4), 1→3 (11), 2→1 (6), 2→3 (2), 3→1 (3).

\[
D^{0} = \begin{pmatrix} 0 & 4 & 11\\ 6 & 0 & 2\\ 3 & \infty & 0 \end{pmatrix}\quad
D^{1} = \begin{pmatrix} 0 & 4 & 11\\ 6 & 0 & 2\\ 3 & \color{fwine}\mathbf{7} & 0 \end{pmatrix}\quad
D^{2} = \begin{pmatrix} 0 & 4 & \color{fwine}\mathbf{6}\\ 6 & 0 & 2\\ 3 & 7 & 0 \end{pmatrix}\quad
D^{3} = \begin{pmatrix} 0 & 4 & 6\\ \color{fwine}\mathbf{5} & 0 & 2\\ 3 & 7 & 0 \end{pmatrix}
\]
\noindent $D^{k}$ allows vertices $1..k$ as intermediates: $d_{ij} \leftarrow \min(d_{ij},\ d_{ik} + d_{kj})$.
Changed entries are in red: $3\to1\to2 = 7$, then $1\to2\to3 = 6$, then $2\to3\to1 = 5$. $D^3$ is final.


Updates: D¹{[}3{]}{[}2{]} = 3 + 4 = 7; D²{[}1{]}{[}3{]} = 4 + 2 = 6; D³{[}2{]}{[}1{]} = 2 + 3 = 5.

\needspace{160pt}

\hypertarget{p2-07--136-01-knapsack-dp-table}{%
\subsection{0/1 Knapsack DP Table}\label{p2-07--136-01-knapsack-dp-table}}

Capacity W = 5; items (weight, value): (2, 3), (3, 4), (4, 5), (5, 6).

\begingroup\small

\par\medskip\noindent\hfil\begin{tabular}{@{}
  >{\raggedright\arraybackslash}p{(\columnwidth - 12\tabcolsep) * \real{0.1067}}
  >{\raggedright\arraybackslash}p{(\columnwidth - 12\tabcolsep) * \real{0.0400}}
  >{\raggedright\arraybackslash}p{(\columnwidth - 12\tabcolsep) * \real{0.0400}}
  >{\raggedright\arraybackslash}p{(\columnwidth - 12\tabcolsep) * \real{0.0400}}
  >{\raggedright\arraybackslash}p{(\columnwidth - 12\tabcolsep) * \real{0.0400}}
  >{\raggedright\arraybackslash}p{(\columnwidth - 12\tabcolsep) * \real{0.0400}}
  >{\raggedright\arraybackslash}p{(\columnwidth - 12\tabcolsep) * \real{0.0400}}@{}}
\toprule\noalign{}
\begin{minipage}[b]{\linewidth}\raggedright
\textbf{i \textbackslash{} w}
\end{minipage} & \begin{minipage}[b]{\linewidth}\raggedright
\textbf{0}
\end{minipage} & \begin{minipage}[b]{\linewidth}\raggedright
\textbf{1}
\end{minipage} & \begin{minipage}[b]{\linewidth}\raggedright
\textbf{2}
\end{minipage} & \begin{minipage}[b]{\linewidth}\raggedright
\textbf{3}
\end{minipage} & \begin{minipage}[b]{\linewidth}\raggedright
\textbf{4}
\end{minipage} & \begin{minipage}[b]{\linewidth}\raggedright
\textbf{5}
\end{minipage} \\
\midrule\noalign{}
0 & 0 & 0 & 0 & 0 & 0 & 0 \\
1 (2, 3) & 0 & 0 & 3 & 3 & 3 & 3 \\
2 (3, 4) & 0 & 0 & 3 & 4 & 4 & \textbf{7} \\
3 (4, 5) & 0 & 0 & 3 & 4 & 5 & 7 \\
4 (5, 6) & 0 & 0 & 3 & 4 & 5 & 7 \\
\bottomrule\noalign{}
\end{tabular}\hfil\null\par\medskip


\endgroup

Answer \textbf{7} (items 1 and 2). Traceback: K{[}4{]}{[}5{]} = K{[}3{]}{[}5{]} = K{[}2{]}{[}5{]} ≠ K{[}1{]}{[}5{]} → item 2 taken; remaining capacity 2 → item 1 taken.

\needspace{95pt}

\hypertarget{p2-07--137-topological-sort-kahns-algorithm}{%
\subsection{Topological Sort (Kahn\textquotesingle s algorithm)}\label{p2-07--137-topological-sort-kahns-algorithm}}

\begin{figure}[!htbp]
\centering
\begin{adjustbox}{max width=\linewidth}
\begin{tikzpicture}[v/.style={vertex, fill=fslate, minimum size=8mm}]
  \node[v, fill=fsage] (A) at (0,1.2) {A}; \node[v, fill=fsage] (B) at (0,-0.4) {B};
  \node[v] (C) at (2.4,1.2) {C}; \node[v] (D) at (2.4,-0.4) {D}; \node[v] (E) at (4.8,0.4) {E};
  \draw[arr] (A) -- (C); \draw[arr] (B) -- (C); \draw[arr] (B) -- (D); \draw[arr] (C) -- (E); \draw[arr] (D) -- (E);
  \foreach \n/\d in {A/0, B/0, C/2, D/1, E/2} \node[smalllbl, text=fwine] at ([yshift=6mm]\n) {in \d};
\end{tikzpicture}
\end{adjustbox}
\caption{A DAG for Kahn's algorithm. Vertices with in-degree 0 (green) start the queue; removing a vertex lowers
the in-degree of its successors.}
\end{figure}


In-degrees: A 0, B 0, C 2, D 1, E 2. Queue A, B → remove A (C: 1) → remove B (C: 0, D: 0) → C → D (E: 0) → E.
One valid order: \textbf{A, B, C, D, E} (B, A, D, C, E is also valid --- topological orders are not unique).

\needspace{191pt}

\hypertarget{p2-07--138-double-hashing}{%
\subsection{Double Hashing}\label{p2-07--138-double-hashing}}

m = 11, h₁(k) = k mod 11, h₂(k) = 7 − (k mod 7), probe i: (h₁ + i·h₂) mod 11. Insert 22, 33 (22 already at 0):
22 → 0. 33 → h₁ = 0 (occupied); h₂ = 7 − 5 = 2 → (0 + 2) mod 11 = \textbf{2}.
h₂ must never be 0 and should be co-prime with m (choose m prime).

\needspace{125pt}

\hypertarget{p2-07--previous-year-questions-pyq-pattern}{%
\section{Previous Year Questions (PYQ pattern)}\label{p2-07--previous-year-questions-pyq-pattern}}

\textbf{Complexity \& recurrences}

\begin{enumerate}
\def\labelenumi{\arabic{enumi}.}
\tightlist
\item
  Solution of T(n) = 2T(n/2) + n log n: \textbf{Θ(n log² n)}
\item
  T(n) = 4T(n/2) + n: \textbf{Θ(n²)}
\item
  T(n) = T(n/2) + n: \textbf{Θ(n)}
\item
  T(n) = 3T(n/4) + n log n: \textbf{Θ(n log n)} (case 3)
\item
  Which is the slowest-growing? (a) n log n (b) n\^{}1.01 (c) n/log n (d) √n log n --- \textbf{(d)}
\item
  Time complexity of the loop \texttt{for(i=1;i\textless{}n;i*=2)}: \textbf{O(log n)}
\item
  Nested \texttt{for(i=1;i\textless{}=n;i++)\ for(j=1;j\textless{}=n;j+=i)}: \textbf{O(n log n)} (harmonic series)
\end{enumerate}

\textbf{Arrays, stacks, queues}

\begin{enumerate}
\def\labelenumi{\arabic{enumi}.}
\setcounter{enumi}{7}
\tightlist
\item
  A{[}1..10{]}{[}1..15{]}, row-major, base 100, 4 bytes/element. Address of A{[}5{]}{[}7{]}: 100 + 4{[}(4)(15) + 6{]} = \textbf{364}
\item
  Postfix of (A + B) * (C − D): \textbf{AB+CD−*}
\item
  Prefix of A + B * C: \textbf{+A*BC}
\item
  Minimum stacks to implement a queue: \textbf{2}
\item
  Circular queue of size n holds at most: \textbf{n − 1} elements (one-slot-empty convention)
\end{enumerate}

\textbf{Trees}

\begin{enumerate}
\def\labelenumi{\arabic{enumi}.}
\setcounter{enumi}{12}
\tightlist
\item
  A binary tree with 20 leaves has how many nodes of degree 2? \textbf{19}
\item
  Maximum nodes in a binary tree of height 5 (root at height 0): 2⁶ − 1 = \textbf{63}
\item
  Number of distinct binary trees with 4 nodes: \textbf{14}
\item
  Inorder: D B E A F C; Preorder: A B D E C F → Postorder: \textbf{D E B F C A}
\item
  Minimum nodes in an AVL tree of height 4: \textbf{12}
\item
  In a complete ternary tree with 10 internal nodes, leaves = 2×10 + 1 = \textbf{21}
\item
  Inorder traversal of BST gives: \textbf{sorted order}
\item
  Worst-case height of BST with n nodes: \textbf{n − 1}
\end{enumerate}

\textbf{Heaps \& hashing}

\begin{enumerate}
\def\labelenumi{\arabic{enumi}.}
\setcounter{enumi}{20}
\tightlist
\item
  Time to build a heap of n elements: \textbf{O(n)}
\item
  Array 89, 19, 50, 17, 12, 15, 2, 5, 7, 11, 6, 9, 100 --- after inserting 100 into max-heap, root = \textbf{100}
\item
  Linear probing suffers from: \textbf{primary clustering}
\item
  With chaining and load factor α, expected unsuccessful search: \textbf{Θ(1 + α)}
\item
  Keys 43, 36, 92, 87, 11, 4, 71, 13, 14 hashed with h(k) = k mod 11 using chaining. Keys in slot 3: \textbf{36 and 14} (43→10, 36→3, 92→4, 87→10, 11→0, 4→4, 71→5, 13→2, 14→3)
\end{enumerate}

\textbf{Sorting \& searching}

\begin{enumerate}
\def\labelenumi{\arabic{enumi}.}
\setcounter{enumi}{25}
\tightlist
\item
  Which sorting algorithm has worst case O(n log n) and is in-place? \textbf{Heap sort}
\item
  Best algorithm for nearly sorted data: \textbf{Insertion sort}
\item
  Quick sort worst case occurs when: \textbf{array is already sorted (pivot = first/last)}
\item
  Which sorts are stable? \textbf{Merge, insertion, bubble, counting, radix}
\item
  Lower bound for comparison sorting: \textbf{Ω(n log n)}
\item
  Maximum comparisons in binary search on 1000 elements: ⌊log₂1000⌋ + 1 = \textbf{10}
\item
  Min and max of n elements need at least: \textbf{⌈3n/2⌉ − 2 comparisons}
\end{enumerate}

\textbf{Design techniques}

\begin{enumerate}
\def\labelenumi{\arabic{enumi}.}
\setcounter{enumi}{32}
\tightlist
\item
  Matrix chain 10×20, 20×30, 30×40: min multiplications = (AB)C = 6000 + 12000 = \textbf{18000}
\item
  LCS of "ABCD" and "ACBD": \textbf{3}
\item
  0/1 knapsack using DP: \textbf{O(nW)} (pseudo-polynomial)
\item
  Huffman coding is an example of: \textbf{greedy}
\item
  Floyd--Warshall uses: \textbf{dynamic programming}
\item
  N-Queens problem is solved by: \textbf{backtracking}
\item
  Which needs both optimal substructure \& overlapping subproblems? \textbf{Dynamic programming}
\end{enumerate}

\textbf{Graphs}

\begin{enumerate}
\def\labelenumi{\arabic{enumi}.}
\setcounter{enumi}{39}
\tightlist
\item
  Dijkstra fails with: \textbf{negative edge weights}
\item
  Bellman-Ford time complexity: \textbf{O(VE)}
\item
  Kruskal\textquotesingle s time complexity: \textbf{O(E log E)}
\item
  BFS is used to find: \textbf{shortest path in unweighted graph}
\item
  Topological sort is possible only for: \textbf{DAG}
\item
  Number of edges in the MST of a connected graph with n vertices: \textbf{n − 1}
\item
  Back edge in DFS of a directed graph indicates: \textbf{a cycle}
\item
  Max-flow equals: \textbf{min-cut capacity}
\end{enumerate}

\textbf{Complexity theory}

\begin{enumerate}
\def\labelenumi{\arabic{enumi}.}
\setcounter{enumi}{47}
\tightlist
\item
  First problem proved NP-complete: \textbf{SAT (Cook--Levin)}
\item
  Which is in P? (a) 3-SAT (b) 2-SAT (c) Clique (d) Hamiltonian cycle --- \textbf{(b)}
\item
  If A ≤ₚ B and B ∈ P then: \textbf{A ∈ P}
\item
  If an NP-complete problem is solved in polynomial time: \textbf{P = NP}
\item
  Halting problem is: \textbf{undecidable, NP-hard but not in NP}
\end{enumerate}

\textbf{Selected topics}

\begin{enumerate}
\def\labelenumi{\arabic{enumi}.}
\setcounter{enumi}{52}
\tightlist
\item
  KMP string matching time: \textbf{O(n + m)}
\item
  FFT multiplies two polynomials of degree n in: \textbf{O(n log n)}
\item
  Randomised quicksort is a: \textbf{Las Vegas algorithm}
\item
  Christofides algorithm gives approximation ratio: \textbf{1.5} (metric TSP)
\end{enumerate}

\textbf{More practice questions}

\begin{enumerate}
\def\labelenumi{\arabic{enumi}.}
\setcounter{enumi}{56}
\tightlist
\item
  After inserting 10, 20, 30 into an empty AVL tree, the root is: \textbf{20}
\item
  Number of rotations needed for the RL case: \textbf{2}
\item
  Build max-heap from {[}4, 10, 3, 5, 1{]}; the resulting array: \textbf{{[}10, 5, 3, 4, 1{]}}
\item
  MST weight for the graph in §13.4: \textbf{11}
\item
  Floyd--Warshall final distance from 2 to 1 in §13.5: \textbf{5}
\item
  0/1 knapsack W = 5 with items (2, 3), (3, 4), (4, 5), (5, 6): \textbf{7}
\item
  Which is NOT a valid topological order of §13.7? (a) A B C D E (b) B A D C E (c) A C B D E (d) B D A C E --- \textbf{Ans: (c)} (C before B violates B → C)
\item
  T(n) = T(n − 1) + n solves to: \textbf{Θ(n²)}
\item
  In double hashing, h₂(k) must never evaluate to: \textbf{0}
\item
  Kruskal\textquotesingle s algorithm uses which data structure to detect cycles? \textbf{Disjoint-set (union--find)}
\end{enumerate}

\begin{revbox}

\begin{itemize}
\tightlist
\item
  Master theorem: compare f(n) with n\^{}(log\_b a)
\item
  n₀ = n₂ + 1 · Catalan 1, 1, 2, 5, 14, 42 · AVL min nodes 1, 2, 4, 7, 12, 20
\item
  Build-heap O(n) · Comparison sort Ω(n log n)
\item
  Dijkstra no negatives · Bellman--Ford O(VE) · Floyd O(V³)
\item
  SAT first NPC · 2-SAT in P · KMP O(n + m)
\end{itemize}

\end{revbox}
